% Part: model-theory
% Chapter: interpolation
% Section: introduction

\documentclass[../../../include/open-logic-section]{subfiles}

\begin{document}

\olfileid{mod}{int}{int}

\section{પરિચય}

અંતર્વેશન પ્રમેય નીચેનું પરિણામ છે: ધારો કે
$\Entails !A \lif !B$. તો એવું !!a{sentence}~$!C$ છે કે
$\Entails !A \lif !C$ અને $\Entails !C \lif !B$. વળી, $!C$માં આવતો
દરેક !!{constant}, !!{function} અને !!{predicate} ($\eq$ સિવાયનો)
$!A$ અને~$!B$ બંનેમાં આવે છે. !!{sentence}~$!C$ને $!A$ અને~$!B$નો
\emph{અંતર્વેશક} કહે છે.

અંતર્વેશન પ્રમેય પોતે જ રસપ્રદ છે, પણ તેનું મુખ્ય મહત્વ એમાં છે કે
તેનો ઉપયોગ કોઈ સિદ્ધાંતમાં વ્યાખ્યેયતા વિશેનાં પરિણામો અને કઈ
શરતો હેઠળ બે સુસંગત સિદ્ધાંતોને જોડવાથી સુસંગત સિદ્ધાંત મળે છે તે
સાબિત કરવા માટે થઈ શકે છે. પહેલા પરિણામને બેથનું વ્યાખ્યેયતા પ્રમેય
અને બીજાને રૉબિન્સનનું સંયુક્ત સુસંગતતા પ્રમેય કહે છે.

\end{document}
